% !PRINTED 642
\weberpage{642}

4. \gesperrt{Every integral rational functional can be decomposed into prime factors, and indeed, apart from unit factors, in only one way.}

\gesperrt{If a product of integral functionals in $\overline{Z}$ is divisible by a prime functional, then at least one of the factors is divisible by this prime functional.}

\begin{center}
§ 178.\\
\gesperrt{Functionals of the field $\overline{\Omega}$.}
\end{center}

A functional of the field $\overline{\Omega}$ is an expression of the form:
\weq{1}{\overline{\omega} = \overline{x}_0 + \overline{x}_1\theta + \cdots + \overline{x}_{n-1}\theta^{n-1},}
where $\theta$ is the algebraic function that generates the field $\Omega$, and $\overline{x}_0, \overline{x}_1, \ldots, \overline{x}_{n-1}$ are \gesperrt{rational functionals}. $\overline{\omega}$ is $= 0$ only when all the coefficients $\overline{x}_0, \overline{x}_1, \ldots, \overline{x}_{n-1}$ vanish.

Every functional which, with respect to $\theta$, is of higher than the $(n-1)$th degree can be replaced by the remainder of the division by $f(\theta)$, and so reduced to the $(n-1)$th degree; and the quotient of two expressions (1), in which the denominator is not $=0$, can be brought into the form (1) as in § 171, 1.
% [sic] printed "in der der Nenner", the relative pronoun should be "dem" (antecedent "der Quotient")

If one takes an arbitrary basis
\[
\omega_1, \omega_2, \ldots, \omega_n
\]
of the field $\Omega$, then $\overline{\omega}$ can be put into the form
\weq{2}{\overline{\omega} = \overline{x}_1\omega_1 + \overline{x}_2\omega_2 + \cdots + \overline{x}_n\omega_n}
as well. If one applies this to the products $\overline{\omega}\omega_i$ and sets
\weq{3}{\overline{\omega}\omega_i = \overline{x}_{i,1}\omega_1 + \overline{x}_{i,2}\omega_2 + \cdots + \overline{x}_{i,n}\omega_n,}
in which the $\overline{x}_{i,n}$ are rational functionals, then there results for $\overline{\omega}$ an equation:
\weq{4}{\left|\begin{array}{lll}
\overline{x}_{1,1} - \overline{\omega}, & \overline{x}_{1,2}, & \ldots, \overline{x}_{1,n} \\
\overline{x}_{2,1}, & \overline{x}_{2,2} - \overline{\omega}, & \ldots, \overline{x}_{2,n} \\
\cdot & \cdot & \cdots \\
\overline{x}_{n,1}, & \overline{x}_{n,2}, & \ldots, \overline{x}_{n,n} - \overline{\omega}
\end{array}\right| = 0,}
and so:

5. \gesperrt{Every functional in $\overline{\Omega}$ satisfies an equation}
\weq{5}{\varphi(\overline{\omega}) = 0,}
\gesperrt{in which}
\weq{6}{\varphi(t) = t^n + \overline{A}_1 t_2^{n-1} + \overline{A}_2 t^{n-2}, \cdots + \overline{A}_n}
% [sic] printed "\overline{A}_1 t_2^{n-1}" and "t^{n-2}, \cdots", misprints for "\overline{A}_1 t^{n-1}" (stray subscript 2 on t) and "t^{n-2} + \cdots" (comma for plus, cf. (1)-(3))
%CONT>>
